\chapter{Finite Fields}

\begin{center}\itshape
At its highest error-correction level, a QR code can restore roughly $30\%$ of its codewords.\\
How can a message repair itself?\\
The answer is polynomial arithmetic in a field with $256$ elements.
\end{center}

In the field $\mathbb F_2=\{0,1\}$, adding $1$ to $1$ gives $0$. There are only two elements. Could there be a field with four elements, or eight, or nine?

A first attempt with four elements is to take the classes of $0,1,X,1+X$ modulo $X^2+X+1$. In this ring,
\[
X^2=X+1
\]
because $-1=1$ in characteristic $2$. The polynomial $X^2+X+1$ has no root in $\mathbb F_2$, so it is irreducible; the quotient is a field. We have built $\mathbb F_4$ using nothing but polynomial arithmetic.

This example raises three questions. Why must the size of every finite field be a prime power? Why is there, up to isomorphism, exactly one field of each such size? And why do its non-zero elements always form a cyclic multiplicative group?

The Frobenius map $x\mapsto x^p$ supplies a large part of the answer. It preserves addition because the intermediate binomial coefficients vanish in characteristic $p$. Its fixed points and iterates reveal the structure of all finite fields.

There is a practical reason to care. Send the values of a low-degree polynomial at several field elements rather than sending only its coefficients. If a few values are damaged, enough others still determine the polynomial. That is the idea behind Reed--Solomon error correction. We will derive the distance bound, then see how a linear system recovers the message.

\section{The characteristic and the Frobenius map}

\begin{proposition}\label{prop:primefield}
Let $K$ be a finite field. There is a prime $p$, the \emph{characteristic} of $K$, such that $p\cdot1=0$ in $K$; the subring generated by $1$ is a field isomorphic to $\mathbb F_p=\Z/p\Z$, the \emph{prime field} of $K$; and $|K|=p^n$, where $n=[K:\mathbb F_p]$.
\end{proposition}
\begin{proof}
The map $\Z\to K$, $k\mapsto k\cdot1$, is a ring homomorphism. $K$ is finite and $\Z$ is not, so its kernel is a non-zero ideal $(p)$, $p\ge1$, and its image is isomorphic to $\Z/p\Z$ (Proposition~\ref{prop:quotring}). The image is a subring of a field, so it is a domain, and $p$ is prime (Proposition~\ref{prop:primemax}: $(p)$ is a prime ideal; $p\ne1$ because $1\ne0$ in $K$). So the image is the field $\mathbb F_p$, and $K$ is a vector space of some finite dimension $n$ over it; choosing a basis gives a bijection $K\cong\mathbb F_p^n$.
\end{proof}

\begin{lemma}[Frobenius]\label{lem:frob}
Let $K$ be a field of characteristic $p$. Then $(a+b)^p=a^p+b^p$ for all $a,b\in K$, and $\varphi:x\mapsto x^p$ is an injective ring homomorphism $K\to K$ that fixes $\mathbb F_p$. If $K$ is finite, $\varphi$ is an automorphism of $K$, the \emph{Frobenius automorphism}.
\end{lemma}
\begin{proof}
By the binomial formula, $(a+b)^p=\sum_k\binom pka^kb^{p-k}$, and for $0<k<p$ the prime $p$ divides $\binom pk=\frac{p!}{k!\,(p-k)!}$ (it divides the numerator, not the denominator). So these terms vanish in characteristic $p$. Products and $1$ are clearly respected. A ring homomorphism out of a field is injective (its kernel is an ideal not containing $1$). On $\mathbb F_p$, $x^p=x$ by Fermat's little theorem. An injective map from a finite set to itself is bijective.
\end{proof}

\begin{example}\label{ex:f4}
$X^2+X+1$ has no root in $\mathbb F_2$, so it is irreducible (degree $2$), and $\mathbb F_4=\mathbb F_2[X]/(X^2+X+1)$ is a field with $4$ elements $0,1,\alpha,\alpha+1$, where $\alpha^2=\alpha+1$. Then $\alpha^3=\alpha^2+\alpha=1$: the group $\mathbb F_4^*=\{1,\alpha,\alpha^2\}$ is cyclic of order $3$. The ring $\Z/4\Z$ also has four elements, but it is not a field ($2\cdot2=0$); a field with $p^n$ elements is \emph{not} $\Z/p^n\Z$ for $n\ge2$.
\end{example}

Every finite field has a multiplicative group of non-zero elements. Rather than examine its multiplication table, we can use the fact that a polynomial of degree $d$ has at most $d$ roots. That simple bound forces the whole group to be cyclic.

\section{The multiplicative group is cyclic}

\begin{theorem}\label{thm:cyclicmult}
Let $K$ be a field. Every finite subgroup $G$ of $K^*$ is cyclic. In particular, if $K$ is finite with $q$ elements, $K^*$ is cyclic of order $q-1$.
\end{theorem}
\begin{proof}
$G$ is a finite abelian group, of order $m$ say, so $G\cong\Z/d_1\oplus\dots\oplus\Z/d_r$ with $d_1\mid\dots\mid d_r$ (Corollary~\ref{cor:abelian}). Every element $x$ of $\bigoplus\Z/d_i$ satisfies $d_rx=0$, so every $x\in G$ satisfies $x^{d_r}=1$. The polynomial $X^{d_r}-1$ has at most $d_r$ roots in the field $K$, so $m\le d_r$. As $m=d_1\cdots d_r$, $r=1$ and $G\cong\Z/m$.
\end{proof}

A generator of $\mathbb F_q^*$ is a \emph{primitive element}. If $\gamma$ is one, every non-zero element is a power of $\gamma$, so $\mathbb F_q=\mathbb F_p(\gamma)\cong\mathbb F_p[X]/(m_\gamma)$ (Theorem~\ref{thm:simple}): every finite field is obtained by Kronecker's construction. Multiplication in $\mathbb F_q^*$ is addition of exponents modulo $q-1$, as on the circle of Figure~\ref{fig:f16}; addition is addition of coordinates in $\mathbb F_p^n$. A finite field is the meeting point of these two very different pictures.

\begin{figure}[ht]
\centering
\fig{ch09-f16}
\caption{Left: $\mathbb F_{16}=\mathbb F_2[\alpha]$ with $\alpha^4=\alpha+1$. Its $15$ non-zero elements are the powers of $\alpha$, placed on a circle, so multiplying turns the circle; under each power are its coordinates on $\alpha^3,\alpha^2,\alpha,1$, which are added bit by bit (exclusive or). Right: the number $N_2(n)$ of irreducible polynomials of degree $n$ over $\mathbb F_2$ follows $2^n/n$ closely (Theorem~\ref{thm:countirr}); the vertical scale is logarithmic.}
\label{fig:f16}
\end{figure}

We know the size of a finite field must be a prime power, and we know its non-zero elements form a cyclic group. Now we must show that a field of each prime-power size exists and that two fields of the same size cannot have different multiplication laws up to isomorphism.

\section{Existence and uniqueness}

\begin{lemma}\label{lem:xqx}
If $K$ is a field with $q$ elements, then $x^q=x$ for every $x\in K$, and
\[
X^q-X=\prod_{a\in K}(X-a)\quad\text{in }K[X].
\]
\end{lemma}
\begin{proof}
For $x\ne0$, $x^{q-1}=1$ by Lagrange's theorem in the group $K^*$ of order $q-1$; multiply by $x$. So the $q$ elements of $K$ are roots of the monic polynomial $X^q-X$ of degree $q$, which is therefore their product $\prod(X-a)$ (each root splits off a factor, Volume~1).
\end{proof}

\begin{theorem}[Galois, Moore]\label{thm:fq}
For every prime $p$ and $n\ge1$ there is a field with $q=p^n$ elements, and any two such fields are isomorphic: it is the splitting field of $X^q-X$ over $\mathbb F_p$. We denote it by $\mathbb F_q$.
\end{theorem}
\begin{proof}
We must show both that a field of each required size exists and that two fields of the same size cannot have different multiplication laws.

\emph{Existence.} Fix $q=p^n$ and begin with the polynomial
\[
F(X)=X^q-X\in\mathbb F_p[X].
\]
Let $L$ be a splitting field of $F$ over $\mathbb F_p$, which exists by Proposition~\ref{prop:splitexists}. Thus $F$ splits completely into linear factors in $L$.

How many different roots does it have? Its formal derivative is
\[
F'(X)=qX^{q-1}-1=-1
\]
because $q=p^n$ vanishes in characteristic $p$. A repeated root of $F$ would also be a root of $F'$, which is impossible. Therefore $F$ has exactly $q$ \emph{distinct} roots in $L$.

Call their set $E$:
\[
E=\{x\in L:x^q=x\}.
\]
We claim that $E$ is a field. The Frobenius law in characteristic $p$ gives
\[
(x+y)^q=x^q+y^q,\qquad (xy)^q=x^qy^q.
\]
Hence the sum and product of two roots are roots again. The same argument with $-y$ gives closure under subtraction, and if $x\ne0$ satisfies $x^q=x$, then
\[
(x^{-1})^q=(x^q)^{-1}=x^{-1},
\]
so every non-zero root has its inverse in $E$. Finally $0,1\in E$. This proves all the field axioms by inheritance from $L$.

Every element of $\mathbb F_p$ satisfies $x^p=x$, and therefore $x^q=x$, so $\mathbb F_p\subseteq E$. Moreover, $E$ contains every root of $F$. Since $L$ is \emph{generated over $\mathbb F_p$ by those roots}, we have $L\subseteq E$, while $E\subseteq L$ by construction. Thus $E=L$ and
\[
|L|=|E|=q.
\]

\emph{Uniqueness.} Suppose $K$ is any field with $q=p^n$ elements. By Proposition~\ref{prop:primefield}, its characteristic is $p$, and it contains a prime subfield identified with $\mathbb F_p$. Every non-zero element of the multiplicative group $K^*$ has order dividing $q-1$, by Lagrange. Hence
\[
x^{q-1}=1\quad(x\in K^*),\qquad x^q=x\quad(x\in K).
\]
This is also Lemma~\ref{lem:xqx}. In other words, all $q$ elements of $K$ are distinct roots of $X^q-X$, and the polynomial splits there. These roots generate all of $K$ over its prime subfield, because they \emph{are} all the elements of $K$. Therefore $K$ is a splitting field of the same polynomial over $\mathbb F_p$.

By the uniqueness theorem for splitting fields (Theorem~\ref{thm:splitunique}), the two fields are isomorphic. We may consequently write $\mathbb F_q$ without ambiguity up to isomorphism.
\end{proof}

\paragraph{Why this works.}
The familiar identity $x^p=x$ holds for the $p$ elements of $\mathbb F_p$. To build a field with $p^n$ elements, we ask for all solutions of $x^{p^n}=x$. Two facts make this work. Its derivative is $-1$, so the polynomial has exactly $p^n$ distinct roots in its splitting field. And in characteristic $p$ the map $x\mapsto x^{p^n}$ preserves addition as well as multiplication, so the roots are closed under the field operations. Thus the roots form a field with exactly the desired size. Uniqueness comes from the reverse observation: every element of an existing field of this size solves the same polynomial.

The equation $x^{p^n}=x$ describes every element of $\mathbb F_{p^n}$. Applying the Frobenius map repeatedly reveals which elements lie in smaller fields, and which automorphisms can fix the prime field.

\section{Subfields and the Galois group}

We use twice that if $d\mid n$, then $a^d-1$ divides $a^n-1$, for integers and for polynomials alike: with $n=de$, $a^n-1=(a^d-1)(a^{d(e-1)}+\dots+a^d+1)$.

\begin{theorem}\label{thm:subfields}
The subfields of $\mathbb F_{p^n}$ are the sets $\{x:x^{p^d}=x\}$ for $d\mid n$; there is exactly one subfield with $p^d$ elements for each divisor $d$ of $n$, and no other.
\end{theorem}
\begin{proof}
Let $K$ be a subfield. It contains the prime field, so $|K|=p^d$ with $d=[K:\mathbb F_p]$, and the tower law gives $n=d\,[\mathbb F_{p^n}:K]$, so $d\mid n$. By Lemma~\ref{lem:xqx}, $K$ is contained in the set of roots of $X^{p^d}-X$, which has at most $p^d$ elements, so $K$ equals that set.

Conversely let $d\mid n$. Then $p^d-1\mid p^n-1$, hence $X^{p^d-1}-1$ divides $X^{p^n-1}-1$, and $X^{p^d}-X$ divides $X^{p^n}-X$. The latter splits in $\mathbb F_{p^n}$ with distinct roots (Lemma~\ref{lem:xqx}), so the former has $p^d$ distinct roots there. They are the fixed points of the homomorphism $\varphi^d$, so they form a subfield, with $p^d$ elements.
\end{proof}

So the subfields of $\mathbb F_{p^n}$ form a copy of the lattice of divisors of $n$: $\mathbb F_{2^{12}}$ contains $\mathbb F_{2^4}$ and $\mathbb F_{2^6}$, which meet in $\mathbb F_{2^2}$, but $\mathbb F_8\not\subseteq\mathbb F_{16}$.

\begin{theorem}\label{thm:galfq}
The automorphisms of $\mathbb F_{p^n}$ form a cyclic group of order $n$, generated by the Frobenius automorphism $\varphi$.
\end{theorem}
\begin{proof}
For $1\le k<n$, $\varphi^k\ne\mathrm{id}$: the polynomial $X^{p^k}-X$ has at most $p^k<p^n$ roots, so some $x$ has $x^{p^k}\ne x$. And $\varphi^n=\mathrm{id}$ by Lemma~\ref{lem:xqx}. So $\langle\varphi\rangle$ has order $n$. Every automorphism fixes $1$, hence the prime field. $\mathbb F_{p^n}$ is a splitting field of $X^{p^n}-X$ over $\mathbb F_p$, a polynomial with distinct roots, so by Theorem~\ref{thm:splitunique} it has exactly $[\mathbb F_{p^n}:\mathbb F_p]=n$ automorphisms. They are the powers of $\varphi$.
\end{proof}

This is the first complete example of the Galois correspondence of Chapter~10: the subgroups of $\Z/n\Z$ (one for each $d\mid n$) match the subfields (one for each $d\mid n$), the subgroup $\langle\varphi^d\rangle$ fixing exactly $\mathbb F_{p^d}$.

Every element of $\mathbb F_{q^n}$ is a root of $X^{q^n}-X$. Factoring this polynomial over $\mathbb F_q$ lets us count irreducible polynomials by degree, just as the divisors of an integer help us count numbers with a given order.

\section{Counting irreducible polynomials}

\begin{think}
How many monic irreducible polynomials of degree $2$ are there over $\mathbb F_2$? Over $\mathbb F_3$? (A monic quadratic is reducible iff it has a root.) Can you guess a formula over $\mathbb F_q$?
\end{think}

Everything above works over any finite field $\mathbb F_q$ in place of $\mathbb F_p$, as $\mathbb F_q\subseteq\mathbb F_{q^n}$ by Theorem~\ref{thm:subfields} (if $q=p^m$, then $m\mid mn$), and $x\mapsto x^q=\varphi^m(x)$ is a homomorphism fixing exactly $\mathbb F_q$.

\begin{theorem}\label{thm:xqnx}
In $\mathbb F_q[X]$, $X^{q^n}-X$ is the product of all monic irreducible polynomials whose degree divides $n$.
\end{theorem}
\begin{proof}
The derivative of $X^{q^n}-X$ is $-1$, so it has no repeated factor: if $P^2$ divided it, every root of $P$ in a splitting field would be a double root. It is monic, and $\mathbb F_q[X]$ is a UFD, so it suffices to show that a monic irreducible $P$, of degree $d$, divides $X^{q^n}-X$ iff $d\mid n$. Let $K=\mathbb F_q[X]/(P)$, a field with $q^d$ elements, and $\alpha$ the class of $X$, whose minimal polynomial is $P$. For $k\ge0$ let $\psi_k(x)=x^{q^k}$, a homomorphism of $K$, and $S=\{k\ge0:\psi_k=\mathrm{id}_K\}$. By Lemma~\ref{lem:xqx}, $d\in S$.

If $d\mid n$, then $\psi_n=(\psi_d)^{n/d}=\mathrm{id}$, so $\alpha^{q^n}=\alpha$, and $P=m_\alpha$ divides $X^{q^n}-X$.

Conversely, suppose $P\mid X^{q^n}-X$, so $\psi_n(\alpha)=\alpha$. The fixed points of $\psi_n$ form a subfield containing $\mathbb F_q$ and $\alpha$, which generate $K$; so $\psi_n=\mathrm{id}$ and $n\in S$. If $k\ge l$ are in $S$, then $\psi_k=\psi_{k-l}\circ\psi_l=\psi_{k-l}$, so $k-l\in S$. Running Euclid's algorithm by subtractions, $g=\gcd(n,d)\in S$: every one of the $q^d$ elements of $K$ is a root of $X^{q^g}-X$, so $q^d\le q^g$, $d\le g$, and $d=g$ divides $n$.
\end{proof}

\begin{theorem}[Gauss]\label{thm:countirr}
Let $N_q(n)$ be the number of monic irreducible polynomials of degree $n$ over $\mathbb F_q$. Then
\[
\sum_{d\mid n}d\,N_q(d)=q^n,\qquad N_q(n)=\frac1n\sum_{d\mid n}\mu(d)\,q^{n/d},
\]
and
\[
\Bigl|N_q(n)-\frac{q^n}n\Bigr|<\frac{2q^{n/2}}n.
\]
In particular $N_q(n)\ge1$ for every $n\ge1$.
\end{theorem}
\begin{proof}
Compare degrees in Theorem~\ref{thm:xqnx}: the irreducibles of degree $d$, for each $d\mid n$, contribute $d\,N_q(d)$. Möbius inversion (Volume~2, Chapter~6), applied to $f(d)=dN_q(d)$ and $g(n)=q^n$, gives the formula. In it, the term $d=1$ is $q^n$, and the others have $n/d\le n/2$, so with $m=\lfloor n/2\rfloor$,
\[
\begin{aligned}\Bigl|nN_q(n)-q^n\Bigr|&\le\sum_{k=1}^mq^k=\frac{q(q^m-1)}{q-1}\eqbr<\frac q{q-1}q^m\le2q^{n/2}.\end{aligned}
\]
For the last claim: $N_q(1)=q$, and for $n\ge2$, $2q^{n/2}\le q^{n/2+1}\le q^n$, so $nN_q(n)>0$.
\end{proof}

So there are $(q^2-q)/2$ irreducible monic quadratics over $\mathbb F_q$: one over $\mathbb F_2$ ($X^2+X+1$) and three over $\mathbb F_3$ ($X^2+1$, $X^2+X+2$, $X^2+2X+2$). Directly: there are $q^2$ monic quadratics, of which $q$ are squares $(X-a)^2$ and $\binom q2$ are products $(X-a)(X-b)$ with $a\ne b$; $q^2-q-\frac{q(q-1)}2=\frac{q^2-q}2$.

\begin{remark}
Think of $X^{q^n}$, the number of monic polynomials of degree $n$, as ``$x$''. Then $n=\log_qx$, and Theorem~\ref{thm:countirr} says that the number of irreducibles of degree $n$ is $\frac x{\log_qx}$ up to an error of size $\sqrt x$: a prime number theorem for $\mathbb F_q[X]$, with an error term as small as the one the Riemann hypothesis would give for the primes (up to a logarithm). For polynomials the proof took a page; for primes it is the most famous open problem in mathematics, met again later in the series. Weil's theorem on curves over finite fields is a deep extension of this ``Riemann hypothesis over finite fields''.
\end{remark}

A finite field is a place where we can add, multiply and divide without leaving a finite alphabet. Evaluating a polynomial at several distinct field elements turns those operations into a way of detecting and correcting errors.

\section{Correcting errors: Reed--Solomon codes}

To send $k$ symbols $m_0,\dots,m_{k-1}\in\mathbb F_q$ over a noisy channel, encode them as the polynomial $f=\sum m_jX^j$ of degree $<k$, choose $n$ distinct points $a_1,\dots,a_n\in\mathbb F_q$ ($k\le n\le q$), and send the \emph{codeword} $(f(a_1),\dots,f(a_n))$. The \emph{Hamming distance} between two words of $\mathbb F_q^n$ is the number of positions where they differ; it satisfies the triangle inequality, since if $u_i\ne w_i$ then $u_i\ne v_i$ or $v_i\ne w_i$.

\begin{proposition}\label{prop:rsdist}
Two different codewords are at distance at least $n-k+1$. Consequently, if at most $t=\lfloor\frac{n-k}2\rfloor$ symbols are corrupted, the sent codeword is the unique codeword closest to the received word.
\end{proposition}
\begin{proof}
If $f\ne g$ have degree $<k$, then $f-g$ is a non-zero polynomial of degree $<k$, so it has at most $k-1$ roots: $f(a_i)=g(a_i)$ for at most $k-1$ indices $i$. If the sent word $c$ and the received word $y$ are at distance $e\le t$, and $c'\ne c$ is another codeword, then $d(c',y)\ge d(c,c')-d(c,y)\ge n-k+1-t>t\ge e$, because $2t\le n-k$.
\end{proof}

Finding the closest codeword by trying all $q^k$ messages is hopeless when $q=256$ and $k=200$. Berlekamp and Welch found a way with linear algebra: an unknown \emph{error locator} polynomial $E$ vanishes at the bad positions, and $fE$ is then determined by a linear system.

\begin{theorem}[Berlekamp--Welch decoding]\label{thm:bw}
Let $y\in\mathbb F_q^n$ differ from the codeword of $f$ in at most $t=\lfloor\frac{n-k}2\rfloor$ positions. Then the linear system
\[
Q(a_i)=y_i\,E(a_i)\qquad(1\le i\le n),
\]
in the unknowns $E$ monic of degree $t$ and $Q$ of degree $<k+t$, has a solution, and every solution satisfies $Q=fE$. So $f$ is the quotient of $Q$ by $E$.
\end{theorem}
\begin{proof}
We must prove two things: there really is a solution of the displayed linear system, and no solution can lead us to the wrong message.

\emph{A solution exists.} Let $B$ be the set of corrupted positions, of size $e\le t$. Choose the monic polynomial
\[
E(X)=X^{t-e}\prod_{i\in B}(X-a_i).
\]
Its degree is exactly $t$. It vanishes at every corrupted evaluation point. Put $Q=fE$, whose degree is less than $k+t$. At a good position, $y_i=f(a_i)$, so $Q(a_i)=y_iE(a_i)$. At a bad position, $E(a_i)=0$, so both sides of the same equation are zero. Thus $(E,Q)$ solves all $n$ equations.

\emph{Every solution gives the sent message.} Now let $(E,Q)$ be \emph{any} solution having the required degrees. We may not assume that its polynomial $E$ has the same roots as the one just constructed. Instead consider
\[
R=Q-fE.
\]
Its degree, unless it is zero, is less than $k+t$. At each of the at least $n-t$ good positions,
\[
R(a_i)=Q(a_i)-f(a_i)E(a_i)
      =y_iE(a_i)-f(a_i)E(a_i)=0.
\]
The evaluation points are distinct, and
\[
n-t\ge k+t
\]
because $2t\le n-k$. So a non-zero polynomial of degree less than $k+t$ would have at least $k+t$ distinct roots, which is impossible. Hence $R=0$, that is, $Q=fE$.

The polynomial $E$ is monic and therefore non-zero. Divide $Q$ by $E$ in $\mathbb F_q[X]$ to recover $f$. This is the unique message compatible with the stated error bound.
\end{proof}

\paragraph{Why this works.}
The product $E(x)$ hides the corrupted positions by making them zeros. At a good position, $y_i=f(a_i)$, so the equation $Q(a_i)=y_iE(a_i)$ says $Q(a_i)=f(a_i)E(a_i)$. At a bad position, both sides vanish because $E(a_i)=0$. We solve for $Q$ and $E$ together using linear equations, even though the message polynomial $f$ is unknown. Why is the result trustworthy? For any candidate solution, $Q-fE$ vanishes at at least $n-t$ points and has degree less than $k+t$. The correction bound $2t\le n-k$ makes those points too numerous for a non-zero polynomial. Therefore every solution gives the same message $f=Q/E$.

The system has $n$ equations and $k+2t\le n$ unknowns (the $t$ lower coefficients of $E$ and the $k+t$ coefficients of $Q$); Gaussian elimination solves it in $O(n^3)$ operations in $\mathbb F_q$, and faster algorithms (Berlekamp--Massey, Euclid's algorithm on polynomials) need $O(n^2)$. QR codes, CDs, DVDs, RAID~6 disk arrays and the Voyager probes all use Reed--Solomon codes over $\mathbb F_{256}$, whose elements are bytes; the highest QR level adds enough redundancy to recover about $30\%$ of the data. The same field, built as $\mathbb F_2[X]/(X^8+X^4+X^3+X+1)$, is the arithmetic of the AES cipher.

\begin{beauty}
A finite field is completely described by two numbers, $p$ and $n$, and yet it carries a rich structure: a cyclic multiplicative group, a cyclic Galois group generated by $x\mapsto x^p$, a lattice of subfields copying the divisors of $n$, and a count of irreducible polynomials that mirrors the distribution of primes. Galois discovered these ``imaginary'' fields in 1830, the year of his work on equations. Today they are everywhere in computing: error-correcting codes, AES, elliptic-curve cryptography, hashing, pseudo-random generators. The Frobenius automorphism, in turn, becomes in Volume~13 the bridge between primes and Galois groups (Frobenius elements, and Chebotarev's theorem).
\end{beauty}

\section{Exercises}

\begin{exercise}[$\star$]\label{ex:f4tables}
Write the addition and multiplication tables of $\mathbb F_4$ (Example~\ref{ex:f4}). Show that $\mathbb F_4$ and $\Z/4\Z$ are not even isomorphic as additive groups, and that $\mathbb F_2[X]/(X^2+1)$ is not a field.
\end{exercise}

\begin{exercise}[$\star$]\label{ex:irrF2}
List the monic irreducible polynomials of degree $\le4$ over $\mathbb F_2$, check the count of Theorem~\ref{thm:countirr}, and factor $X^{16}-X$ in $\mathbb F_2[X]$.
\end{exercise}

\begin{exercise}[$\star\star$]\label{ex:primpoly}
An irreducible $P\in\mathbb F_p[X]$ is \emph{primitive} if the class of $X$ generates $(\mathbb F_p[X]/(P))^*$. Show that $X^4+X+1$ is primitive over $\mathbb F_2$, but $X^4+X^3+X^2+X+1$ is not. How many primitive elements does $\mathbb F_{16}$ have, and how many primitive polynomials of degree $4$ are there over $\mathbb F_2$?
\end{exercise}

\begin{exercise}[$\star\star$]\label{ex:minus1}
Let $q$ be odd. Show that $-1$ is a square in $\mathbb F_q$ iff $q\equiv1\pmod4$. Deduce that $-1$ is a square in $\mathbb F_9$, and find a square root of $-1$ in $\mathbb F_9=\mathbb F_3[X]/(X^2+X+2)$.
\end{exercise}

\begin{exercise}[$\star\star$]\label{ex:twosq}
Show that every element of a finite field is a sum of two squares.
\end{exercise}

\begin{exercise}[$\star\star$]\label{ex:powersums}
\begin{enumerate}[label=(\alph*)]
\item Show that for an integer $k\ge0$, $\displaystyle\sum_{x\in\mathbb F_q}x^k$ is $-1$ if $k\ge1$ and $(q-1)\mid k$, and $0$ otherwise (with $0^0=1$).
\item (Wilson) Show that the product of all non-zero elements of $\mathbb F_q$ is $-1$.
\end{enumerate}
\end{exercise}

\begin{exercise}[$\star\star$]\label{ex:rs7}
Over $\mathbb F_7$, a message $f$ of degree $<2$ is sent as $(f(1),f(2),\dots,f(6))$, and $y=(5,3,2,4,0,1)$ is received. Assuming at most two errors, decode it with Theorem~\ref{thm:bw}.
\end{exercise}

\begin{exercise}[$\star\star\star$]\label{ex:chevalley}
(Chevalley--Warning.) Let $F\in\mathbb F_q[X_1,\dots,X_n]$ have total degree $d<n$. Show that the number $N$ of zeros of $F$ in $\mathbb F_q^n$ is divisible by $p$, the characteristic. Deduce that if $F(0)=0$, then $F$ has a non-zero zero; for instance, every quadratic form in $3$ variables over $\mathbb F_q$ has a non-trivial zero.
\end{exercise}

\begin{exercise}[$\star\star$]\label{pb:fibmodp}
Let $p\ne2,5$ be a prime and $\left(\frac5p\right)$ the Legendre symbol. Let $(F_n)$ be the Fibonacci numbers.
\begin{enumerate}[label=(\alph*)]
\item Show that $p$ divides $F_{p-(5/p)}$: so $p\mid F_{p-1}$ when $p\equiv\pm1\pmod5$, and $p\mid F_{p+1}$ when
$p\equiv\pm2\pmod5$.
\item Show that the period of $(F_n\bmod p)$ divides $p-1$ in the first case and $2(p+1)$ in the second.
\item Find the periods for $p=7$ and $p=11$.
\end{enumerate}
\end{exercise}

\begin{exercise}[$\star$]\label{pb:pascal2}
Let $s(n)$ be the number of $1$s in the binary expansion of $n$.
\begin{enumerate}[label=(\alph*)]
\item Show that $\binom{2^k+m}j\equiv\binom mj+\binom m{j-2^k}\pmod2$ for $0\le m<2^k$ (with $\binom mi=0$ for $i<0$),
and explain the pattern of Figure~\ref{fig:pascal}.
\item Deduce that row $n$ of Pascal's triangle contains exactly $2^{s(n)}$ odd numbers.
\item Show that the rows $0,1,\dots,2^k-1$ contain $3^k$ odd numbers in all, and that the proportion of odd entries in them
tends to $0$.
\end{enumerate}
\end{exercise}

\begin{exercise}[$\star\star\star$]\label{pb:lucaslehmer}
\emph{The Lucas--Lehmer test.} Let $s_0=4$ and $s_{k+1}=s_k^2-2$. Let $p$ be an odd prime and $M_p=2^p-1$. Show that $M_p$
is prime if and only if $M_p\mid s_{p-2}$.
\end{exercise}

\begin{exercise}[$\star\star$]\label{pb:moore}
\emph{The Moore determinant.} Let $q$ be a prime power and $\Delta(x_1,\dots,x_n)=\det\bigl(x_i^{q^{j-1}}\bigr)_{1\le i,j\le n}$.
\begin{enumerate}[label=(\alph*)]
\item (Moore, 1896.) Show that
\[
\Delta(x_1,\dots,x_n)=\prod_{i=1}^n\ \prod_{c\in\mathbb F_q^{i-1}}\bigl(x_i+c_1x_1+\dots+c_{i-1}x_{i-1}\bigr)
\]
in $\mathbb F_q[x_1,\dots,x_n]$.
\item Deduce that $\alpha_1,\dots,\alpha_n\in\mathbb F_{q^m}$ are linearly independent over $\mathbb F_q$ if and only if
$\Delta(\alpha_1,\dots,\alpha_n)\ne0$.
\item Let $V\subset\mathbb F_{q^m}$ be an $\mathbb F_q$-subspace. Show that $P_V(X)=\prod_{v\in V}(X-v)$ is a linear combination
of $X,X^q,X^{q^2},\dots$, so that $P_V(x+y)=P_V(x)+P_V(y)$.
\end{enumerate}
\end{exercise}

\begin{exercise}[$\star\star$]\label{pb:mseq}
\emph{Maximal length sequences.} Let $\alpha$ generate the cyclic group $\mathbb F_{2^n}^*$, $f=x^n+c_{n-1}x^{n-1}+\dots+c_0$ its
minimal polynomial over $\mathbb F_2$, $\operatorname{Tr}(\beta)=\beta+\beta^2+\dots+\beta^{2^{n-1}}$, and
$s_k=\operatorname{Tr}(\alpha^k)\in\mathbb F_2$.
\begin{enumerate}[label=(\alph*)]
\item Show that $s_{k+n}=c_{n-1}s_{k+n-1}+\dots+c_0s_k$ for all $k$.
\item Show that the map $\beta\mapsto\bigl(\operatorname{Tr}(\beta),\operatorname{Tr}(\beta\alpha),\dots,\operatorname{Tr}(\beta\alpha^{n-1})\bigr)$ is a bijection
$\mathbb F_{2^n}\to\mathbb F_2^n$. Deduce that $(s_k)$ has period $2^n-1$, that every non-zero block of $n$ bits appears exactly
once among the windows $(s_k,\dots,s_{k+n-1})$, $0\le k<2^n-1$, and that every non-zero sequence satisfying the recurrence of (a)
is a shift of $(s_k)$.
\item Show that a period contains $2^{n-1}$ ones and $2^{n-1}-1$ zeros, and that for $0<t<2^n-1$,
$\sum_{k=0}^{2^n-2}(-1)^{s_k+s_{k+t}}=-1$.
\end{enumerate}
\end{exercise}

\begin{figure}[ht]
\centering
\fig{ch01-pascal}
\caption{The odd entries of the first $32$ rows of Pascal's triangle. Rows $16$ to $31$ repeat rows $0$ to $15$ twice, side
by side (Problem~\ref{pb:pascal2}).}
\label{fig:pascal}
\end{figure}

\section{Hints}

\begin{hint}{ex:f4tables}
In $\mathbb F_4$, $x+x=0$ for every $x$. In $\mathbb F_2[X]$, $X^2+1=(X+1)^2$.
\end{hint}
\begin{hint}{ex:irrF2}
A polynomial of degree $2$ or $3$ is irreducible iff it has no root; one of degree $4$ iff it has no root and is not $(X^2+X+1)^2$.
\end{hint}
\begin{hint}{ex:primpoly}
The order of $X$ divides $15$; compute $X^3$ and $X^5$. For the counts, each primitive polynomial of degree $4$ has $4$ roots, all primitive.
\end{hint}
\begin{hint}{ex:minus1}
$-1$ is the unique element of order $2$ in the cyclic group $\mathbb F_q^*$. Try elements $a+b\alpha$ with $\alpha^2=-\alpha-2$.
\end{hint}
\begin{hint}{ex:twosq}
For $q$ odd, how many squares are there? Count the values of $x^2$ and of $a-y^2$.
\end{hint}
\begin{hint}{ex:powersums}
(a) Write the non-zero elements as powers of a primitive element: a geometric sum. (b) $X^{q-1}-1=\prod_{a\ne0}(X-a)$; look at the constant term.
\end{hint}
\begin{hint}{ex:rs7}
Here $n=6$, $k=2$, $t=2$: unknowns $E=X^2+e_1X+e_0$ and $Q$ of degree $\le3$. If you guess where the errors are, the existence part of the proof tells you $E$ and $Q$.
\end{hint}
\begin{hint}{ex:chevalley}
In $\mathbb F_q$, $1-F(x)^{q-1}$ is $1$ if $F(x)=0$ and $0$ otherwise. Sum over $x\in\mathbb F_q^n$, expand $F^{q-1}$ into monomials and use Exercise~\ref{ex:powersums}(a).
\end{hint}

\begin{hint}{pb:fibmodp}
Binet's formula holds in any field where $X^2-X-1$ has two distinct roots; the Frobenius $x\mapsto x^p$ permutes them.
\end{hint}

\begin{hint}{pb:pascal2}
$(1+x)^{2^k}\equiv1+x^{2^k}$ in $\mathbb F_2[x]$.
\end{hint}

\begin{hint}{pb:lucaslehmer}
$s_k=\omega^{2^k}+\bar\omega^{2^k}$ with $\omega=2+\sqrt3$, $\omega\bar\omega=1$. One direction: the order of $\omega$ in
$(\Z/q\Z)[\sqrt3]$. The other: $\omega=\frac{(1+\sqrt3)^2}2$ and Frobenius in $\mathbb F_{M_p^2}$.
\end{hint}

\begin{hint}{pb:moore}
(a) Expand along the last row, and note that the last row is a combination of the others when $x_n$ is an $\mathbb F_q$-combination
of $x_1,\dots,x_{n-1}$.
\end{hint}

\begin{hint}{pb:mseq}
(b) The trace is a non-zero linear form, and $(\beta,\gamma)\mapsto\operatorname{Tr}(\beta\gamma)$ is non-degenerate.
(c) $s_k+s_{k+t}=\operatorname{Tr}(\alpha^k(1+\alpha^t))$.
\end{hint}

\section{Solutions}

\begin{solution}{ex:f4tables}
\idea Compute with $\alpha^2=\alpha+1$ and $2=0$; compare the orders of elements in the two additive groups.
\writeup Writing $\beta=\alpha+1=\alpha^2$: addition is $x+x=0$, $1+\alpha=\beta$, $1+\beta=\alpha$, $\alpha+\beta=1$ (and $0$ is neutral); multiplication is $\alpha\cdot\alpha=\beta$, $\alpha\beta=\alpha^3=1$, $\beta\beta=\alpha^4=\alpha$ (and $1$ is neutral, $0$ absorbing). In $(\mathbb F_4,+)$ every element satisfies $2x=0$, so it is $(\Z/2\Z)^2$; in $\Z/4\Z$, $1$ has order $4$. So they are not isomorphic groups, and in fact the additive group of any field of characteristic $p$ has exponent $p$. Finally $X^2+1=(X+1)^2$ in $\mathbb F_2[X]$, so the class of $X+1$ is a non-zero element of $\mathbb F_2[X]/(X^2+1)$ whose square is $0$: the ring is not a domain.
\end{solution}

\begin{solution}{ex:irrF2}
\idea For degrees $2$ and $3$, a polynomial is irreducible exactly when it has no root in $\mathbb F_2$. In degree $4$ we must also rule out a factor of degree $2$. There is only one monic irreducible quadratic to test. Check the resulting list against the counting formula before forming the degree-$16$ product.

\writeup Over $\mathbb F_2$ a polynomial has a root at $0$ precisely when its constant coefficient vanishes, and a root at $1$ precisely when the number of non-zero coefficients is even.

The irreducible polynomials of degree $1$ are
\[
X,\qquad X+1.
\]
The only monic quadratic with neither root is $X^2+X+1$, so this is the sole irreducible quadratic. A cubic is reducible if and only if it has a root, which gives exactly
\[
X^3+X+1,\qquad X^3+X^2+1
\]
as irreducible cubics.

For degree $4$, the four monic polynomials with no root in $\mathbb F_2$ are
\[
\begin{gathered}
F_1=X^4+X+1,\qquad F_2=X^4+X^3+1,\\
F_3=X^4+X^3+X^2+X+1,\qquad
F_4=X^4+X^2+1.
\end{gathered}
\]
A root-free quartic is reducible only if it factors into irreducible quadratics; the only candidate factor is $P=X^2+X+1$. In the quotient by $P$, let $\alpha=X\bmod P$. Then $\alpha^2=\alpha+1$, $\alpha^3=1$, and $\alpha^4=\alpha$. We obtain
\[
F_1(\alpha)=1,\quad
F_2(\alpha)=\alpha,\quad
F_3(\alpha)=\alpha+1,\quad
F_4(\alpha)=0.
\]
In fact $F_4=P^2$ in characteristic $2$. Thus the irreducible quartics are precisely $F_1,F_2,F_3$.

The counts in degrees $1,2,3,4$ are consequently $2,1,2,3$. The formula
\[
N_2(d)=\frac1d\sum_{e\mid d}\mu(e)\,2^{d/e}
\]
gives, independently,
\[
N_2(1)=2,\quad N_2(2)=\frac{4-2}{2}=1,\quad
N_2(3)=\frac{8-2}{3}=2,\quad
N_2(4)=\frac{16-4}{4}=3.
\]

Finally, Theorem~\ref{thm:xqnx} says that $X^{16}-X$ is the product of \emph{all} monic irreducibles whose degrees divide $4$, each appearing once. Degree $3$ therefore does not occur. The factorisation is
\[
\begin{aligned}
X^{16}-X={}&X(X+1)(X^2+X+1)\\
&{}\cdot(X^4+X+1)(X^4+X^3+1)\\
&{}\cdot(X^4+X^3+X^2+X+1).
\end{aligned}
\]
It is square-free since its formal derivative is $-1=1$ in characteristic $2$. The total degree is $1+1+2+4+4+4=16$, as expected.
\end{solution}

\begin{solution}{ex:primpoly}
\idea There are $16-1=15$ non-zero elements of $\mathbb F_{16}$. An element is primitive if its multiplicative order is $15$. Since the prime divisors of $15$ are $3$ and $5$, it is enough to rule out orders dividing $3$ or $5$.

\writeup Let $P=X^4+X+1$. By Exercise~\ref{ex:irrF2}, $P$ is irreducible over $\mathbb F_2$. Thus
\[
K=\mathbb F_2[X]/(P)
\]
is a field with $16$ elements. Write $x$ for the class of $X$. Its minimal polynomial has degree $4$, so $1,x,x^2,x^3$ form an $\mathbb F_2$-basis. In particular $x^3\ne1$.

The relation $P(x)=0$ gives $x^4=x+1$ in characteristic $2$, and hence
\[
x^5=x(x+1)=x^2+x\ne1
\]
by linear independence of $1,x,x^2$. By Lagrange, the order of $x$ divides $15$. Every proper divisor of $15$ is $1$, $3$, or $5$. The two inequalities $x^3\ne1$ and $x^5\ne1$ rule out all proper orders. Therefore $x$ has order $15$, and $P$ is primitive.

Now take $Q=X^4+X^3+X^2+X+1$. It is irreducible by the same previous exercise. In characteristic $2$,
\[
(X+1)Q=X^5+1=X^5-1.
\]
The class $y$ of $X$ modulo $Q$ therefore satisfies $y^5=1$. Since $y\ne1$ (otherwise its degree-four minimal polynomial would have root $1$), its order is exactly $5$, not $15$. So $Q$ is irreducible but \emph{not} primitive.

The group $\mathbb F_{16}^*$ is cyclic of order $15$, by Theorem~\ref{thm:cyclicmult}. The number of its generators is
\[
\varphi(15)=15(1-1/3)(1-1/5)=8.
\]
Every primitive element $\gamma$ has a minimal polynomial of degree $4$ over $\mathbb F_2$: it generates all of the multiplicative group and therefore all of the field. Its four Frobenius conjugates
\[
\gamma,\gamma^2,\gamma^4,\gamma^8
\]
are distinct and primitive; distinctness follows because the orbit under squaring has length $4$, and primitivity follows because each exponent $2^j$ is coprime to $15$.

Those four conjugates have the same degree-four minimal polynomial, and no primitive element outside the orbit has that polynomial. The $8$ primitive elements therefore split into $8/4=2$ disjoint orbits. Hence there are
\[
\boxed{8\text{ primitive elements and }2\text{ primitive quartics}}
\]
over $\mathbb F_2$. They are $X^4+X+1$ and $X^4+X^3+1$.
\end{solution}

\begin{solution}{ex:minus1}
\idea In a cyclic group, a given element is a square iff its ``exponent'' with respect to a generator is even.
\writeup Let $\gamma$ generate $\mathbb F_q^*$, of even order $q-1$. The element $\gamma^{(q-1)/2}$ has order $2$, and it is the only one, as $X^2-1$ has the two roots $\pm1$; so $-1=\gamma^{(q-1)/2}$. The squares in $\mathbb F_q^*$ are the $(\gamma^j)^2=\gamma^{2j}$; as $q-1$ is even, the parity of an exponent $i$ of $\gamma^i$ is well defined modulo $q-1$, so $\gamma^i$ is a square iff $i$ is even. So $-1$ is a square iff $(q-1)/2$ is even, iff $q\equiv1\pmod4$. As $9\equiv1\pmod4$, $-1$ is a square in $\mathbb F_9$. With $\alpha^2=-\alpha-2=2\alpha+1$: $(\alpha+2)^2=\alpha^2+4\alpha+4=\alpha^2+\alpha+1=3\alpha+2=2=-1$. So $\alpha+2$ is a square root of $-1$ (and so is $-(\alpha+2)=2\alpha+1$).
\end{solution}

\begin{solution}{ex:twosq}
\idea The two sets of possible values $x^2$ and $a-y^2$ must overlap if their sizes add to more than the number $q$ of field elements. The size of the set of squares depends on whether the characteristic is $2$ or odd.

\writeup Let $K=\mathbb F_q$ and fix $a\in K$. We seek $x,y\in K$ satisfying
\[
x^2+y^2=a.
\]

If $q$ is even, then $q=2^n$. The Frobenius map $x\mapsto x^2$ is injective because $x^2=y^2$ implies $(x-y)^2=0$, hence $x=y$ in a field. An injective self-map of a finite set is bijective. So there exists $x$ with $x^2=a$; simply choose $y=0$.

Now suppose $q$ is odd. If $x^2=y^2$, then
\[
(x-y)(x+y)=0
\]
forces $x=y$ or $x=-y$. For $x\ne0$, these are two different elements because the characteristic is not $2$. The map $x\mapsto x^2$ is therefore two-to-one on $K^*$, so there are $(q-1)/2$ non-zero squares. Including $0$, the set
\[
S=\{x^2:x\in K\}
\]
has exactly $(q+1)/2$ elements.

Consider also
\[
T=\{a-y^2:y\in K\}=a-S.
\]
Translation by $a$ and negation are bijections, so $|T|=|S|=(q+1)/2$. Both lie inside the $q$-element set $K$. Since
\[
|S|+|T|=q+1>q,
\]
the pigeonhole principle gives $S\cap T\ne\varnothing$. Pick $z\in S\cap T$. By definition, $z=x^2=a-y^2$ for suitable $x,y\in K$, and hence $a=x^2+y^2$.

The argument works for \emph{every} $a$ and \emph{every} finite field, not just prime fields. It proves existence; it does not count the number of representations of each $a$.
\end{solution}

\begin{solution}{ex:powersums}
\idea The non-zero elements form a cyclic group of order $q-1$. Sum their $k$-th powers as a geometric progression. For the product identity, either pair each element with its inverse or compare the constant term of the polynomial having all non-zero field elements as roots.

\writeup Let $K=\mathbb F_q$ and choose a generator $\gamma$ of $K^*$, which has order $q-1$ by Theorem~\ref{thm:cyclicmult}.

\emph{(a) Power sums.} First take $k=0$. With the convention $0^0=1$, every term in
\[
\sum_{x\in K}x^0
\]
equals $1$, so the sum is $q\cdot1=0$ \emph{as an element of $K$}, since the characteristic is $p$ and $q=p^n$ vanishes in $K$.

Now let $k\ge1$. The term from $x=0$ equals $0$. Every non-zero field element is a unique power $\gamma^j$ with $0\le j\le q-2$, so
\[
\sum_{x\in K}x^k
=\sum_{j=0}^{q-2}(\gamma^k)^j.
\]
If $(q-1)\mid k$, then $\gamma^k=1$, every term equals $1$, and the sum is $q-1=-1$ in $K$.

If $(q-1)\nmid k$, then $\gamma^k\ne1$ and the finite geometric-series identity applies:
\[
\sum_{j=0}^{q-2}(\gamma^k)^j
=\frac{(\gamma^k)^{q-1}-1}{\gamma^k-1}
=\frac{1-1}{\gamma^k-1}=0.
\]
The denominator is non-zero precisely because $\gamma^k\ne1$. Hence
\[
\boxed{\sum_{x\in\mathbb F_q}x^k=
\begin{cases}
-1,&k\ge1\text{ and }(q-1)\mid k,\\
0,&\text{otherwise},
\end{cases}}
\]
where both values are interpreted inside the field.

\emph{(b) Wilson's product identity.} By Lemma~\ref{lem:xqx}, the polynomial $X^q-X$ has every element of $K$ as a root, and all are distinct. Dividing out the root $0$ gives
\[
X^{q-1}-1=\prod_{a\in K^*}(X-a).
\]
Compare the constant coefficients:
\[
-1=(-1)^{q-1}\prod_{a\in K^*}a.
\]
If $q$ is odd, then $q-1$ is even, and the product equals $-1$. If $q$ is even, the field has characteristic $2$ and $-1=1$, so the product again equals $-1$. Therefore
\[
\boxed{\prod_{a\in\mathbb F_q^*}a=-1}.
\]
For a prime field, this reduces to Wilson's congruence $(p-1)!\equiv-1\pmod p$. It continues to hold in all finite extension fields, where the expression $(q-1)!$ would not be the right interpretation of the product.
\end{solution}

\begin{solution}{ex:rs7}
\idea Write the six linear equations $Q(a_i)=y_iE(a_i)$ and solve; or guess the error positions and check, since Theorem~\ref{thm:bw} says that any solution gives the right answer.
\writeup The unknowns are $e_0,e_1$ and $Q=q_0+q_1X+q_2X^2+q_3X^3$. Solving the $6\times6$ system modulo $7$ by elimination (the matrix is invertible) gives $E=X^2+3$ and $Q=2X^3+3X^2+6X+2$. Check, with $f=Q/E$: indeed $(2X+3)(X^2+3)=2X^3+3X^2+6X+9$, and $9=2$ in $\mathbb F_7$. The table of values at $a=1,\dots,6$ confirms the six equations:
\begin{center}
\small
\begin{tabular}{c|cccccc}
$a$ & $1$ & $2$ & $3$ & $4$ & $5$ & $6$\\\hline
$y$ & $5$ & $3$ & $2$ & $4$ & $0$ & $1$\\
$E(a)$ & $4$ & $0$ & $5$ & $5$ & $0$ & $4$\\
$yE(a)=Q(a)$ & $6$ & $0$ & $3$ & $6$ & $0$ & $4$\\
$f(a)=2a+3$ & $5$ & $0$ & $2$ & $4$ & $6$ & $1$
\end{tabular}
\end{center}
So the message is $f=3+2X$. The error locator $E=X^2+3=(X-2)(X-5)$ shows where the errors were: positions $2$ and $5$, where $3$ and $0$ were received instead of $0$ and $6$.
\end{solution}

\begin{solution}{ex:chevalley}
\idea Count zeros inside $\mathbb F_q$ with an indicator polynomial, and show that the polynomial is too small in degree to have a non-zero sum.
\writeup By Lemma~\ref{lem:xqx}, $F(x)^{q-1}=1$ if $F(x)\ne0$, so $1-F(x)^{q-1}$ is the indicator of $F(x)=0$, and in $\mathbb F_q$,
\[
\begin{aligned}N\cdot1&=\sum_{x\in\mathbb F_q^n}\bigl(1-F(x)^{q-1}\bigr)\eqbr=q^n\cdot1-\sum_{x}F(x)^{q-1}=-\sum_xF(x)^{q-1}.\end{aligned}
\]
Expand $F^{q-1}$ into monomials $cX_1^{k_1}\cdots X_n^{k_n}$ of total degree $k_1+\dots+k_n\le d(q-1)<n(q-1)$. For each, some $k_i<q-1$, and
\[
\sum_{x\in\mathbb F_q^n}x_1^{k_1}\cdots x_n^{k_n}=\prod_{j=1}^n\Bigl(\sum_{x_j\in\mathbb F_q}x_j^{k_j}\Bigr)
\]
has the factor $\sum_{x_i}x_i^{k_i}=0$ (Exercise~\ref{ex:powersums}(a): $k_i=0$, or $k_i\ge1$ not divisible by $q-1$). So $N\cdot1=0$ in $\mathbb F_q$, i.e.\ $p\mid N$. If $F(0)=0$, then $N\ge1$, so $N\ge p\ge2$, and there is a zero other than $0$. A quadratic form in $3$ variables has degree $2<3$ and vanishes at $0$.
\end{solution}

\begin{solution}{pb:fibmodp}
\idea Binet's formula $F_n=\frac{\varphi^n-\psi^n}{\varphi-\psi}$ is an identity between sequences satisfying the same
recurrence, so it holds in every field where $X^2-X-1$ has two distinct roots, for instance in a field with $p$ or $p^2$
elements. In such a field the Frobenius $x\mapsto x^p$ is an automorphism; it either fixes the roots (when $5$ is a square
modulo $p$) or swaps them. Both cases give a power of $\varphi$ equal to $1$ or to $-1$.

\writeup The discriminant of $X^2-X-1$ is $5\ne0$ in $\mathbb F_p$. If $\left(\frac5p\right)=1$ its roots lie in $K=\mathbb
F_p$; otherwise it is irreducible over $\mathbb F_p$ and its roots lie in $K=\mathbb F_p[X]/(X^2-X-1)$, a field with $p^2$
elements (Volume~3, Chapter~9). Call the roots $\varphi\ne\psi$; then $\varphi+\psi=1$ and $\varphi\psi=-1$. The sequences
$(F_n\bmod p)$ and $\bigl(\frac{\varphi^n-\psi^n}{\varphi-\psi}\bigr)$ both satisfy $u_{n+2}=u_{n+1}+u_n$ in $K$, with
$u_0=0$, $u_1=1$; so they coincide. The map $x\mapsto x^p$ is a field homomorphism of $K$ fixing $\mathbb F_p$ (Volume~3,
Chapter~9), so it maps a root of $X^2-X-1$ to a root.

\emph{Case $\left(\frac5p\right)=1$.} Then $\varphi,\psi\in\mathbb F_p^*$, so $\varphi^{p-1}=\psi^{p-1}=1$ by Fermat. Hence
$F_{n+p-1}=F_n$ in $\mathbb F_p$ for all $n$, and $F_{p-1}=F_0=0$.

\emph{Case $\left(\frac5p\right)=-1$.} Then $\varphi\notin\mathbb F_p$. The elements of $K$ fixed by $x\mapsto x^p$ are the
roots of $X^p-X$, at most $p$ of them, and the $p$ elements of $\mathbb F_p$ are roots; so $\varphi^p\ne\varphi$, and
$\varphi^p=\psi$. Likewise $\psi^p=\varphi$. So $\varphi^{p+1}=\varphi\psi=-1=\psi^{p+1}$, which gives
$F_{p+1}=\frac{-1-(-1)}{\varphi-\psi}=0$, and $\varphi^{2(p+1)}=\psi^{2(p+1)}=1$, so $F_{n+2(p+1)}=F_n$ for all $n$.

This proves (a). For (b): the set of $N\in\Z$ such that $F_{n+N}\equiv F_n$ for all $n\ge0$ (extend the sequence to negative
indices by the recurrence) is a subgroup of $\Z$, generated by the period; it contains $p-1$, respectively $2(p+1)$, so the
period divides it. Finally, by quadratic reciprocity $\left(\frac5p\right)=\left(\frac p5\right)$ (Volume~2, Chapter~5), and the
squares modulo $5$ are $\pm1$; this gives the two cases as stated.

(c) For $p=7\equiv2\pmod5$, the period divides $16$; the sequence modulo $7$ is $0$, $1$, $1$, $2$, $3$, $5$, $1$, $6$, $0$, $6$, $6$, $5$, $4$, $2$, $6$, $1$, $0$, $1$, \dots,
of period $16$. For $p=11\equiv1\pmod5$ it divides $10$: $0$, $1$, $1$, $2$, $3$, $5$, $8$, $2$, $10$, $1$, $0$, $1$, \dots, of period $10$.
\tools Volume~3, Chapter~9 (finite fields, Frobenius); Volume~2, Chapter~5 (quadratic reciprocity); Volume~1, Chapter~8 (linear
recurrences, Binet).
\leads The period $\pi(p)$ is called the Pisano period. Wall asked in 1960 whether $\pi(p^2)\ne\pi(p)$ for every prime $p$; a
prime with $\pi(p^2)=\pi(p)$ would be a \emph{Wall--Sun--Sun prime}, and none is known (they are linked to the first case of
Fermat's last theorem). Direct computation gives $\pi(p)$ for all $p<10^5$.
\end{solution}

\begin{solution}{pb:pascal2}
\idea Row $n$ of Pascal's triangle is the polynomial $(1+x)^n$. Modulo $2$, squaring is additive (Frobenius), so
$(1+x)^{2^k}=1+x^{2^k}$, and multiplying by it copies a row twice, shifted. That is the self-similar picture; the count follows
by peeling off binary digits one at a time.

\writeup (a) In $\mathbb F_2[x]$, $(u+v)^2=u^2+v^2$, so by induction $(1+x)^{2^k}=1+x^{2^k}$ (Volume~3, Chapter~9). Hence, for
$0\le m<2^k$,
\[
(1+x)^{2^k+m}=(1+x^{2^k})(1+x)^m=(1+x)^m+x^{2^k}(1+x)^m\quad\text{in }\mathbb F_2[x],
\]
and comparing the coefficients of $x^j$ gives the claim. As $(1+x)^m$ has degree $m<2^k$, the two copies do not overlap: row
$2^k+m$ is row $m$, then a gap of zeros, then row $m$ again. This is the pattern of Figure~\ref{fig:pascal}: the triangle of
rows $0$ to $2^{k+1}-1$ consists of three copies of the triangle of rows $0$ to $2^k-1$.

(b) Let $N(n)$ be the number of odd entries of row $n$. By (a), $N(2^k+m)=2N(m)$ for $0\le m<2^k$, and $N(0)=1$. Removing the
leading binary digit of $n$ halves $N$ and decreases $s$ by one, so $N(n)=2^{s(n)}$ by induction on $n$. (This is also
Lucas's theorem, Volume~2, Chapter~3: $\binom nj$ is odd iff every binary digit of $j$ is at most the corresponding digit of
$n$, which gives $2^{s(n)}$ choices of $j$.)

(c) Write $n<2^k$ with $k$ binary digits $n_0,\dots,n_{k-1}$. Then $2^{s(n)}=\prod_i2^{n_i}$, and summing over all digit
choices gives $\sum_{n<2^k}2^{s(n)}=\prod_{i<k}(1+2)=3^k$. The rows $0$ to $2^k-1$ contain $\sum_{n<2^k}(n+1)=2^{k-1}(2^k+1)$
entries, so the proportion of odd ones is $\frac{3^k}{2^{k-1}(2^k+1)}\le2\bigl(\frac34\bigr)^k\to0$.
\tools Volume~3, Chapter~9 (Frobenius in characteristic $p$); Volume~2, Chapter~3 (Lucas's theorem); Volume~1, Chapter~2.
\leads The limit shape of Figure~\ref{fig:pascal}, rescaled, is the Sierpiński triangle, of dimension $\frac{\log3}{\log2}$
(Volume~7 defines the dimension). Modulo a prime $p$ the same argument gives $\prod(n_i+1)$ non-zero entries in row $n$ (Fine,
1947).
\end{solution}

\begin{solution}{pb:lucaslehmer}
\idea The recursion $s\mapsto s^2-2$ is the doubling formula for $x+x^{-1}$: if $s=x+x^{-1}$ then $s^2-2=x^2+x^{-2}$. With
$x=\omega=2+\sqrt3$, $s_k=\omega^{2^k}+\omega^{-2^k}$, and $M_p\mid s_{p-2}$ says that $\omega^{2^{p-1}}=-1$, so $\omega$ has order $2^p$,
modulo $M_p$. If $M_p$ had a small prime factor $q$, the group where $\omega$ lives modulo $q$ would be too small for that order.
Conversely, if $M_p$ is prime, the Frobenius of $\mathbb F_{M_p^2}$ computes $\omega^{(M_p+1)/2}$, once $\omega$ is written as a square.

\writeup Let $M=M_p$, $\omega=2+\sqrt3$, $\bar\omega=2-\sqrt3$; $\omega\bar\omega=1$. By induction $s_k=\omega^{2^k}+\bar\omega^{2^k}$:
$s_0=4$ and $(\omega^{2^k}+\bar\omega^{2^k})^2-2=\omega^{2^{k+1}}+\bar\omega^{2^{k+1}}$.

\emph{If $M\mid s_{p-2}$, then $M$ is prime.} Let $q$ be a prime factor of $M$ ($q\ne2$, and $q\ne3$ since $2^p\equiv2\pmod3$ makes
$M\equiv1\pmod3$). Work in the finite ring $R=(\Z/q\Z)[X]/(X^2-3)$, with $q^2$ elements, where $\sqrt3$ is the class of $X$. In $R$,
$\omega^{2^{p-2}}+\bar\omega^{2^{p-2}}=0$; multiply by $\omega^{2^{p-2}}$: $\omega^{2^{p-1}}=-1$, so $\omega^{2^p}=1$ and $\omega$ is a unit of order
$2^p$ (as $-1\ne1$ in $R$). The units of $R$ form a group with at most $q^2-1$ elements, so $2^p\le q^2-1$. If $M$ were composite, its
smallest prime factor would satisfy $q^2\le M=2^p-1$, and $2^p\le2^p-2$: absurd.

\emph{If $M$ is prime, then $M\mid s_{p-2}$.} As $p\ge3$, $M\equiv7\pmod8$, so $\bigl(\frac2M\bigr)=1$. And $M\equiv1\pmod3$, $M\equiv3\pmod4$,
so by reciprocity $\bigl(\frac3M\bigr)=-\bigl(\frac M3\bigr)=-1$ (Volume~2, Chapter~5). So $X^2-3$ is irreducible over $\mathbb F_M$ and
$K=\mathbb F_M(\sqrt3)$ is a field with $M^2$ elements (Volume~3, Chapter~9). Its Frobenius $x\mapsto x^M$ is a ring homomorphism fixing
$\mathbb F_M$, and $\sqrt3^M=3^{(M-1)/2}\sqrt3=-\sqrt3$ by Euler's criterion. Let $\tau=1+\sqrt3$; then $\tau^2=4+2\sqrt3=2\omega$ and
$\tau^{M+1}=\tau\cdot\tau^M=(1+\sqrt3)(1-\sqrt3)=-2$. Hence, with $\frac{M+1}2=2^{p-1}$,
\[
\omega^{2^{p-1}}=\Bigl(\frac{\tau^2}2\Bigr)^{(M+1)/2}=\frac{\tau^{M+1}}{2\cdot2^{(M-1)/2}}=\frac{-2}{2\bigl(\frac2M\bigr)}=-1.
\]
Then $s_{p-2}=\omega^{2^{p-2}}+\omega^{-2^{p-2}}=\omega^{-2^{p-2}}\bigl(\omega^{2^{p-1}}+1\bigr)=0$ in $K$. As $s_{p-2}$ is an integer, $M\mid s_{p-2}$.
\tools Volume~2, Chapter~5 (reciprocity, supplementary laws, Euler's criterion); Volume~3, Chapters~6 and~9 (quotient rings, finite
fields, Frobenius).
\leads Lucas (1876) proved $M_{127}$ prime by hand, a record that stood until computers; D.~H.~Lehmer (1930) gave the clean form.
Every record prime since 1952, except a few, is a Mersenne prime found by this test, now by the GIMPS project; the largest
known has more than $41$ million digits. The test is the analogue, in the group of norm-one elements of $\mathbb F_{M^2}$, whose
order $M+1$ is a power of $2$, of Pépin's test for Fermat numbers.
\end{solution}

\begin{solution}{pb:moore}
\idea The Frobenius $x\mapsto x^q$ is $\mathbb F_q$-linear. So the Moore matrix has dependent rows whenever the $x_i$ are
dependent over $\mathbb F_q$, and this gives all the linear factors of the determinant.

\writeup (a) Induction on $n$; for $n=1$, $\Delta=x_1$. Expanding along the last row, $\Delta(x_1,\dots,x_n)$ is a polynomial in
$x_n$ of degree $q^{n-1}$ with leading coefficient $\Delta_{n-1}=\Delta(x_1,\dots,x_{n-1})$. For $c\in\mathbb F_q^{n-1}$, substitute
$x_n=-\sum c_jx_j$: since $(\sum c_jx_j)^{q^k}=\sum c_jx_j^{q^k}$ (Frobenius is additive and fixes $\mathbb F_q$), the last row becomes
$-\sum c_j(\text{row }j)$, and the determinant vanishes. So each $x_n+\sum c_jx_j$ divides $\Delta$ in
$\mathbb F_q(x_1,\dots,x_{n-1})[x_n]$. These $q^{n-1}$ linear factors are pairwise distinct, so their product, of degree $q^{n-1}$ in
$x_n$, divides $\Delta$, and comparing leading coefficients
\[
\Delta(x_1,\dots,x_n)=\Delta_{n-1}\prod_{c\in\mathbb F_q^{n-1}}\Bigl(x_n+\sum_jc_jx_j\Bigr).
\]
The induction hypothesis on $\Delta_{n-1}$ finishes the proof (the identity holds in $\mathbb F_q[x_1,\dots,x_n]$ since both sides lie
there).

(b) Evaluate (a) at $\alpha$: $\Delta(\alpha)=0$ if and only if $\alpha_i+\sum_{j<i}c_j\alpha_j=0$ for some $i$ and some
$c\in\mathbb F_q^{i-1}$, that is, if and only if the $\alpha_i$ are linearly dependent over $\mathbb F_q$.

(c) Let $\alpha_1,\dots,\alpha_n$ be a basis of $V$. By (a),
\[
\Delta(\alpha_1,\dots,\alpha_n,X)=\Delta(\alpha)\prod_{v\in V}(X+v),
\]
and $\prod_{v\in V}(X+v)=\prod_{v\in V}(X-v)$ since $V=-V$. Expanding the determinant along its last row, $(X,X^q,\dots,X^{q^n})$,
gives a combination of $X,X^q,\dots,X^{q^n}$. Divide by $\Delta(\alpha)\ne0$. Each $X^{q^k}$ is additive, so $P_V$ is too.

\tools Volume~3, Chapter~9 (Frobenius); Volume~1, Chapter~8 (determinants).
\leads Moore's determinant is the $q$-analogue of the Vandermonde determinant. Linearised polynomials built this way give the
Gabidulin codes used in network coding, and Dickson's invariants of $GL_n(\mathbb F_q)$, the coefficients of
$\prod_{v\in\mathbb F_q^n}(X-v)$ written in the $x_i$.
\end{solution}

\begin{solution}{pb:mseq}
\idea $s_k$ is a fixed linear form evaluated at $\alpha^k$, and $\alpha^k$ runs over all of $\mathbb F_{2^n}^*$. Every statistic
of the sequence becomes a statistic of a linear form on the whole field, which is perfectly balanced.

\writeup (a) $\operatorname{Tr}$ is $\mathbb F_2$-linear (Frobenius is additive). From $f(\alpha)=0$, $\alpha^{k+n}=\sum_ic_i\alpha^{k+i}$,
and applying $\operatorname{Tr}$ gives the recurrence (signs do not matter in characteristic $2$).

(b) The trace is a polynomial of degree $2^{n-1}$, so it is not zero on all $2^n$ elements; it is a non-zero linear form onto
$\mathbb F_2$. If $\operatorname{Tr}(\beta\alpha^i)=0$ for $0\le i<n$, then $\operatorname{Tr}(\beta\gamma)=0$ for all $\gamma$, since the
$\alpha^i$ span $\mathbb F_{2^n}$ (it has degree $n$); if $\beta\ne0$, $\beta\gamma$ runs over the whole field, which is impossible.
So the linear map is injective, hence bijective. The window at $k$ is the image of $\beta=\alpha^k$. As $k$ runs over
$0,\dots,2^n-2$, $\alpha^k$ runs over the non-zero elements once each, so the windows are the non-zero $n$-bit blocks, once each.
Distinct windows mean the period is exactly $2^n-1$ (a window determines all later terms by (a), and $\alpha^{2^n-1}=1$). A sequence
satisfying (a) is determined by its first window; if that window is non-zero it equals the window of $(s_k)$ at some $k_0$, and then
the sequence is $(s_{k+k_0})$.

(c) The ones in a period are the non-zero $\beta$ with $\operatorname{Tr}\beta=1$; the kernel of $\operatorname{Tr}$ has $2^{n-1}$
elements, so $2^{n-1}$ have trace $1$, and $2^{n-1}-1$ non-zero ones have trace $0$. For the correlation,
$s_k+s_{k+t}=\operatorname{Tr}(\alpha^k(1+\alpha^t))$ with $1+\alpha^t\ne0$; as $k$ varies, $\alpha^k(1+\alpha^t)$ runs over the
non-zero elements, and
\[
\sum_{\beta\ne0}(-1)^{\operatorname{Tr}\beta}=\sum_\beta(-1)^{\operatorname{Tr}\beta}-1=0-1=-1.
\]

\tools Volume~3, Chapter~9 (the cyclic group $\mathbb F_q^*$, the trace; Exercise~9.3 for primitive polynomials); Volume~1, Chapter~8.
\leads These are Golomb's three randomness postulates (1967): balance, runs, two-valued autocorrelation. For $x^4+x+1$ the period is
$000100110101111$. The sequence is useless for cryptography, since $2n$ bits reveal the recurrence (Berlekamp--Massey), but its
flat correlation is what GPS satellites use: each one transmits a Gold code, the sum of two $m$-sequences of degree $10$.
\end{solution}

\section*{Sources and further reading}

\begin{itemize}
\item \'E.~Galois, \emph{Sur la théorie des nombres}, Bulletin des sciences mathématiques de Férussac 13 (1830), 428--435. The first construction of the fields $\mathbb F_{p^n}$, with ``imaginary roots'' of irreducible congruences.
\item E.~H.~Moore, \emph{A doubly-infinite system of simple groups}, Bull.\ New York Math.\ Soc.\ 3 (1893), 73--78. Every finite field is one of Galois' fields.
\item C.~F.~Gauss, \emph{Disquisitiones generales de congruentiis} (the unpublished Section~VIII of the \emph{Disquisitiones}, printed in his Werke, vol.~2, 1863). The count of irreducible polynomials modulo $p$.
\item I.~S.~Reed and G.~Solomon, \emph{Polynomial codes over certain finite fields}, J.\ SIAM 8 (1960), 300--304. The codes of this chapter, in five pages.
\item L.~R.~Welch and E.~R.~Berlekamp, \emph{Error correction for algebraic block codes}, US Patent 4\,633\,470 (1986). The decoding algorithm of Theorem~\ref{thm:bw}.
\item C.~Chevalley, \emph{Démonstration d'une hypothèse de M.~Artin}, Abh.\ Math.\ Sem.\ Hamburg 11 (1935), 73--75; E.~Warning, \emph{Bemerkung zur vorstehenden Arbeit von Herrn Chevalley}, same volume, 76--83. Exercise~\ref{ex:chevalley}.
\item R.~Lidl and H.~Niederreiter, \emph{Introduction to Finite Fields and Their Applications}, revised ed., Cambridge University Press, 1994. Everything about finite fields, with codes and cryptography.
\item J.~Daemen and V.~Rijmen, \emph{The Design of Rijndael: AES, the Advanced Encryption Standard}, Springer, 2002. Why AES computes in $\mathbb F_{256}$, an application of finite-field arithmetic.
\item S.~B.~Wicker and V.~K.~Bhargava (eds.), \emph{Reed--Solomon Codes and Their Applications}, IEEE Press, 1994. Applications to digital storage and communication; the 1960 Reed--Solomon paper is cited above.
\end{itemize}
